\absheading{Resumen}
\begin{absblock}
El objetivo de la investigación fue construir un prototipo funcional que clasifique el nivel de tueste de granos de café a partir de fotografías con Teachable Machine; el alcance de esta primera etapa se limita a capturar imágenes, entrenar el modelo y verificarlo, y deja la acción física automatizada para una etapa posterior. Se desarrolló un estudio aplicado con enfoque cuantitativo y diseño experimental, con 3.158 fotografías de siete clases (seis estados del grano y una clase de fondo), tomadas en una caja con iluminación controlada sobre dos fondos de papel; la colección se dividió por bloques de cinco fotos consecutivas en 2.208 para entrenamiento, 635 para validación y 315 para prueba, y el modelo se evaluó con las fotos separadas y según el fondo. Los resultados mostraron una exactitud de 99,1\,\% en validación (629 de 635) y de 99,4\,\% en prueba (313 de 315); los ocho errores aparecieron sobre el fondo blanco, con 97,2\,\% de aciertos frente a 100\,\% sobre el fondo anterior. En conclusión, el prototipo clasificó el tueste con alta exactitud en condiciones controladas, aunque el fondo influyó en el resultado, por lo que cada clase debe compartir el mismo fondo y la misma luz.

\textbf{Palabras clave:} café, tueste, visión artificial, clasificación, Teachable Machine
\end{absblock}

\absheading{Abstract}
\begin{absblock}
The objective of the research was to build a functional prototype that classifies the roast level of coffee beans from photographs using Teachable Machine; the scope of this first stage is limited to capturing images, training the model, and verifying it, leaving the automated physical action for a later stage. An applied study with a quantitative approach and an experimental design was conducted, using 3,158 photographs of seven classes (six bean states and one background class) taken in a box with controlled lighting on two paper backgrounds; the collection was split into blocks of five consecutive photos: 2,208 for training, 635 for validation, and 315 for testing, and the model was evaluated on the held-out photos and by background. The results showed an accuracy of 99.1\,\% in validation (629 of 635) and 99.4\,\% in testing (313 of 315); all eight errors occurred on the white background, with 97.2\,\% correct versus 100\,\% on the earlier background. In conclusion, the prototype classified roast level with high accuracy under controlled conditions, although the background influenced the result, so every class must share the same background and lighting.

\textbf{Keywords:} coffee, roasting, computer vision, classification, Teachable Machine
\end{absblock}

\absheading{Resumo}
\begin{absblock}
O objetivo da pesquisa foi construir um protótipo funcional que classifique o nível de torra de grãos de café a partir de fotografias com o Teachable Machine; o escopo desta primeira etapa limita-se a capturar imagens, treinar o modelo e verificá-lo, deixando a ação física automatizada para uma etapa posterior. Foi desenvolvido um estudo aplicado com abordagem quantitativa e delineamento experimental, com 3.158 fotografias de sete classes (seis estados do grão e uma classe de fundo), tiradas em uma caixa com iluminação controlada sobre dois fundos de papel; a coleção foi dividida em blocos de cinco fotos consecutivas: 2.208 para treinamento, 635 para validação e 315 para teste, e o modelo foi avaliado com as fotos separadas e conforme o fundo. Os resultados mostraram exatidão de 99,1\,\% na validação (629 de 635) e de 99,4\,\% no teste (313 de 315); os oito erros ocorreram sobre o fundo branco, com 97,2\,\% de acertos frente a 100\,\% sobre o fundo anterior. Em conclusão, o protótipo classificou o nível de torra com alta exatidão em condições controladas, embora o fundo tenha influenciado o resultado, de modo que cada classe deve compartilhar o mesmo fundo e a mesma luz.

\textbf{Palavras-chave:} café, torra, visão computacional, classificação, Teachable Machine
\end{absblock}
